% ============================================================
% fig_graphical_model.tex
% Modelo gráfico generativo do modelo fatorial
% Requer no preâmbulo: \usepackage{tikz}
%   \usetikzlibrary{arrows.meta, positioning, calc}
% ============================================================

\begin{figure}[htbp]
\centering
\begin{tikzpicture}[
    font=\small,
    % nó de variável latente (círculo cinza claro)
    latent/.style={
        circle,
        draw=gray!60,
        fill=gray!12,
        line width=0.7pt,
        minimum size=1.05cm,
        inner sep=2pt
    },
    % nó de variável observada (círculo duplo / mais escuro)
    observed/.style={
        circle,
        draw=gray!70,
        fill=gray!22,
        line width=0.7pt,
        minimum size=1.05cm,
        inner sep=2pt,
        double,
        double distance=1.5pt
    },
    % setas
    arr/.style={-{Stealth[length=5pt,width=4pt]}, gray!70, line width=0.7pt},
    % seta de loading (mais destacada)
    loading/.style={-{Stealth[length=5pt,width=4pt]}, gray!55, line width=0.5pt},
    % label de parâmetro nas setas
    edgelabel/.style={font=\footnotesize, fill=white, inner sep=1.5pt,
                      midway, sloped}
]

% ── Fatores latentes (linha superior) ───────────────────────
\node[latent] (f1) at (0, 0)   {$f_1$};
\node[latent] (f2) at (3, 0)   {$f_2$};
\node[font=\small, gray] at (5.2, 0) {$\cdots$};
\node[latent] (fk) at (6.6, 0) {$f_K$};

% Label da camada
\node[font=\small\itshape, gray!70] at (-1.8, 0) {fatores};

% ── Ativos observados (linha inferior) ──────────────────────
\node[observed] (r1) at (-0.5, -3.0) {$r_1$};
\node[observed] (r2) at ( 2.0, -3.0) {$r_2$};
\node[observed] (r3) at ( 4.5, -3.0) {$r_3$};
\node[font=\small, gray] at (5.9, -3.0) {$\cdots$};
\node[observed] (rn) at ( 7.1, -3.0) {$r_N$};

% Label da camada
\node[font=\small\itshape, gray!70] at (-1.8, -3.0) {retornos};

% ── Ruído idiossincrático (linha inferior, abaixo dos ativos) ─
\node[latent, minimum size=0.85cm] (e1) at (-0.5, -4.8) {$\varepsilon_1$};
\node[latent, minimum size=0.85cm] (e2) at ( 2.0, -4.8) {$\varepsilon_2$};
\node[latent, minimum size=0.85cm] (e3) at ( 4.5, -4.8) {$\varepsilon_3$};
\node[latent, minimum size=0.85cm] (en) at ( 7.1, -4.8) {$\varepsilon_N$};

% ── Setas fatores → ativos ──────────────────────────────────
% f1 → r1 (carregamento forte)
\draw[loading, line width=1.1pt]
    (f1) -- node[edgelabel, xshift=-3pt]{$\beta_{11}$} (r1);

% f1 → r2
\draw[loading, line width=0.9pt]
    (f1) -- node[edgelabel]{$\beta_{21}$} (r2);

% f1 → r3 (carregamento fraco / tracejado)
\draw[loading, dashed, opacity=0.4]
    (f1) -- (r3);

% f2 → r1 (fraco)
\draw[loading, dashed, opacity=0.4]
    (f2) -- (r1);

% f2 → r2
\draw[loading, line width=0.9pt]
    (f2) -- node[edgelabel]{$\beta_{22}$} (r2);

% f2 → r3 (forte)
\draw[loading, line width=1.1pt]
    (f2) -- node[edgelabel]{$\beta_{32}$} (r3);

% fk → rn
\draw[loading, line width=1.1pt]
    (fk) -- node[edgelabel]{$\beta_{NK}$} (rn);

% fk → r3 (fraco)
\draw[loading, dashed, opacity=0.4]
    (fk) -- (r3);

% ── Setas ε → ativos ────────────────────────────────────────
\draw[arr] (e1) -- (r1);
\draw[arr] (e2) -- (r2);
\draw[arr] (e3) -- (r3);
\draw[arr] (en) -- (rn);

% ── Legenda ─────────────────────────────────────────────────
\node[font=\footnotesize, gray!80, align=left] at (3.3, -5.9) {%
    \raisebox{2pt}{\tikz{\draw[loading, line width=1.1pt, -stealth](0,0)--(0.7,0);}}~
    exposição fatorial $\beta_{ik}$ \quad
    \raisebox{2pt}{\tikz{\draw[loading, dashed, opacity=0.5, -stealth](0,0)--(0.7,0);}}~
    exposição nula (esparsidade) \quad
    \raisebox{2pt}{\tikz{\draw[arr, -stealth](0,0)--(0.7,0);}}~
    ruído idiossincrático $\varepsilon_i$%
};

\end{tikzpicture}
\caption{Interpretação generativa do modelo fatorial: um conjunto de $K$ fatores
latentes $f_k$ propaga-se aos retornos dos $N$ ativos por meio das exposições
fatoriais $\beta_{ik}$, enquanto o componente idiossincrático $\varepsilon_i$
captura a variação não explicada pelos fatores comuns e não se propaga entre
ativos. Setas tracejadas indicam exposições nulas, ilustrando a esparsidade
natural de modelos com fatores de setor e país \parencite{paleologo2025}.}
\label{fig:graphical_model}
\end{figure}